\section{Interacting Fields}

\subsection{Types of Interaction (Common Examples)}
Interactions are added to the free Lagrangian:
\begin{equation}
\mathcal{L}=\mathcal{L}_0+\mathcal{L}_{\mathrm{int}}.
\end{equation}

\paragraph{Scalar interactions.}
\begin{equation}
\mathcal{L}_{\mathrm{int}}=-\frac{\lambda}{4!}\phi^4,
\qquad
\mathcal{L}_{\mathrm{int}}=-\frac{g}{3!}\phi^3.
\end{equation}

\paragraph{Yukawa interaction (scalar + fermion).}
\begin{equation}
\mathcal{L}_{\mathrm{int}}=-g\,\bar{\psi}\psi\,\phi.
\end{equation}

\paragraph{QED interaction.}
\begin{equation}
\mathcal{L}_{\mathrm{int}}=-e\,\bar{\psi}\gamma^\mu\psi\,A_\mu.
\end{equation}

\paragraph{Key point.}
The interaction is treated perturbatively when the coupling is small.

\subsection{The Interaction Picture}
Split the Hamiltonian:
\begin{equation}
H = H_0 + H_I,
\end{equation}
where $H_0$ is exactly solvable (free theory), $H_I$ is the interaction.

Interaction-picture operators evolve with $H_0$:
\begin{equation}
\mathcal{O}_I(t)=e^{iH_0 t}\mathcal{O}_S e^{-iH_0 t}.
\end{equation}
States evolve with $H_I$:
\begin{equation}
i\frac{d}{dt}|\psi(t)\rangle_I = H_I(t)|\psi(t)\rangle_I,
\qquad
H_I(t)=e^{iH_0 t}H_I e^{-iH_0 t}.
\end{equation}

Define the time-evolution operator in the interaction picture:
\begin{equation}
|\psi(t)\rangle_I = U_I(t,t_0)|\psi(t_0)\rangle_I,
\qquad
i\frac{d}{dt}U_I(t,t_0)=H_I(t)U_I(t,t_0),
\quad U_I(t_0,t_0)=\mathbf{1}.
\end{equation}

\subsection{Dyson's Formula (Dyson Series)}
Formal solution:
\begin{equation}
U_I(t,t_0)=T\exp\left(-i\int_{t_0}^{t} dt'\,H_I(t')\right).
\end{equation}
Expanding:
\begin{align}
U_I(t,t_0)
&=\mathbf{1}
-i\int_{t_0}^{t} dt_1\,H_I(t_1)
+\frac{(-i)^2}{2!}\int_{t_0}^{t} dt_1dt_2\,T\{H_I(t_1)H_I(t_2)\}
+\cdots
\end{align}

In covariant form, for local interactions:
\begin{equation}
S \equiv U_I(\infty,-\infty)
=
T\exp\left(i\int d^4x\,\mathcal{L}_{\mathrm{int}}(x)\right).
\end{equation}

\subsection{Scattering and the $S$-Matrix}
Define asymptotic in/out states (free particle states):
\begin{equation}
|i\rangle = |\text{in}\rangle,
\qquad
|f\rangle = |\text{out}\rangle.
\end{equation}
The $S$-matrix maps in-states to out-states:
\begin{equation}
|f\rangle = S|i\rangle.
\end{equation}
Matrix element:
\begin{equation}
S_{fi}=\langle f|S|i\rangle.
\end{equation}

It is conventional to write
\begin{equation}
S = \mathbf{1}+iT,
\end{equation}
so that the scattering amplitude is encoded in $T$.

\paragraph{Momentum conservation.}
For translational invariance,
\begin{equation}
\langle f|iT|i\rangle
=
i(2\pi)^4\delta^{(4)}(p_f-p_i)\,\mathcal{M}_{fi},
\end{equation}
where $\mathcal{M}$ is the invariant amplitude.

\subsection{Wick's Theorem}
For free fields, time-ordered products can be expanded into normal-ordered products plus contractions.

For scalar fields:
\begin{equation}
T\{\phi(x)\phi(y)\}
=
:\phi(x)\phi(y):
+\contraction{}{\phi(x)}{}{\phi(y)}\phi(x)\phi(y),
\end{equation}
where the contraction is defined as
\begin{equation}
\contraction{}{\phi(x)}{}{\phi(y)}\phi(x)\phi(y)
\equiv
\langle 0|T\{\phi(x)\phi(y)\}|0\rangle
= \Delta_F(x-y).
\end{equation}

General statement:
\begin{equation}
T\{\phi_1\phi_2\cdots\phi_n\}
=
:\phi_1\phi_2\cdots\phi_n:
+\text{(all possible contractions)}.
\end{equation}

\subsection{Feynman Diagrams (Meaning)}
In perturbation theory:
\begin{itemize}
\item Expand $S$ in powers of the coupling.
\item Use Wick's theorem to evaluate vacuum expectation values.
\item Each term corresponds to a sum of diagrams.
\end{itemize}
Diagrams provide a bookkeeping method for contractions.

\subsection{Feynman Rules (Scalar $\phi^4$ Theory)}
Take
\begin{equation}
\mathcal{L}=\frac12(\partial\phi)^2-\frac12 m^2\phi^2-\frac{\lambda}{4!}\phi^4.
\end{equation}
Momentum-space Feynman rules:
\begin{itemize}
\item External line: attach momentum $p$ (on-shell $p^2=m^2$).
\item Propagator (internal line with momentum $k$):
\begin{equation}
\frac{i}{k^2-m^2+i\epsilon}.
\end{equation}
\item Vertex ($\phi^4$):
\begin{equation}
-i\lambda.
\end{equation}
\item Integrate over undetermined loop momenta:
\begin{equation}
\int \frac{d^4k}{(2\pi)^4}.
\end{equation}
\item Impose momentum conservation at each vertex:
\begin{equation}
(2\pi)^4\delta^{(4)}\!\left(\sum p_{\mathrm{in}}-\sum p_{\mathrm{out}}\right).
\end{equation}
\item Divide by symmetry factors when needed.
\end{itemize}

\subsection{Amplitudes and the Invariant Matrix Element}
For $2\to 2$ scattering:
\begin{equation}
\langle p_3,p_4|iT|p_1,p_2\rangle
=
i(2\pi)^4\delta^{(4)}(p_1+p_2-p_3-p_4)\,\mathcal{M}.
\end{equation}

\paragraph{Tree-level example in $\phi^4$.}
At lowest order,
\begin{equation}
\mathcal{M}_{\text{tree}} = -\lambda.
\end{equation}

\subsection{Decay Rates and Cross Sections (Core Formulas)}
\paragraph{Lorentz-invariant phase space.}
For a final state of $n$ particles:
\begin{equation}
d\Pi_n = (2\pi)^4\delta^{(4)}\!\left(p_i-\sum_{j=1}^n p_j\right)
\prod_{j=1}^n \frac{d^3p_j}{(2\pi)^3\,2E_{p_j}}.
\end{equation}

\paragraph{Decay rate.}
For $A\to 1+\cdots+n$:
\begin{equation}
d\Gamma = \frac{1}{2m_A}\,|\mathcal{M}|^2\, d\Pi_n.
\end{equation}

\paragraph{Cross section.}
For $2\to n$:
\begin{equation}
d\sigma = \frac{1}{4\sqrt{(p_1\cdot p_2)^2-m_1^2m_2^2}}\,|\mathcal{M}|^2\, d\Pi_n.
\end{equation}
In the center-of-mass (COM) frame, the flux factor simplifies to
\begin{equation}
4E_1E_2|\mathbf{v}_1-\mathbf{v}_2|.
\end{equation}

\subsection{Green's Functions (Correlation Functions)}
Define the $n$-point time-ordered Green function:
\begin{equation}
G^{(n)}(x_1,\dots,x_n)
=
\langle 0|T\{\phi(x_1)\cdots\phi(x_n)\}|0\rangle.
\end{equation}
These are the central objects in perturbation theory.

\subsection{Connected Diagrams and Vacuum Bubbles}
The vacuum-to-vacuum amplitude:
\begin{equation}
\langle 0|S|0\rangle
=
\langle 0|T\exp\left(i\int d^4x\,\mathcal{L}_{\mathrm{int}}(x)\right)|0\rangle.
\end{equation}
Diagrams with no external legs are \textbf{vacuum bubbles}. They factorize and cancel in
properly normalized correlation functions:
\begin{equation}
\langle \Omega|T\{\phi(x_1)\cdots\phi(x_n)\}|\Omega\rangle
=
\frac{\langle 0|T\{\phi(x_1)\cdots\phi(x_n)\,S\}|0\rangle}{\langle 0|S|0\rangle}.
\end{equation}

\paragraph{Connected vs disconnected.}
The generating functional (in path integral language) satisfies:
\begin{equation}
Z[J]=e^{iW[J]},
\end{equation}
where $W[J]$ generates \textbf{connected} diagrams.
(You may be expected to know this identity even in a canonical approach.)

\subsection{Reduction Formula (LSZ)}
The LSZ reduction formula relates scattering amplitudes to time-ordered correlators.

For a real scalar field, schematically:
\begin{equation}
\langle p_1',\dots,p_m'|S|p_1,\dots,p_n\rangle
\;\propto\;
\prod_{i=1}^m \left[\int d^4x_i\, e^{ip_i'\cdot x_i}(\square_{x_i}+m^2)\right]
\prod_{j=1}^n \left[\int d^4y_j\, e^{-ip_j\cdot y_j}(\square_{y_j}+m^2)\right]
\langle 0|T\{\phi(x_1)\cdots\phi(y_n)\}|0\rangle.
\end{equation}
Interpretation:
\begin{itemize}
\item Apply $(\square+m^2)$ to \emph{amputate} external propagators.
\item Fourier transform to put external legs on-shell.
\item The result gives the $S$-matrix element (up to normalization factors).
\end{itemize}

\subsection{Summary (Exam Checklist)}
\begin{itemize}
\item Interaction picture: operators evolve with $H_0$, states with $H_I$.
\item Dyson formula:
\[
S=T\exp\left(i\int d^4x\,\mathcal{L}_{\mathrm{int}}(x)\right).
\]
\item Wick theorem reduces time-ordered products to contractions $\Delta_F$.
\item Feynman rules: propagator $i/(p^2-m^2+i\epsilon)$, vertices from $\mathcal{L}_{\mathrm{int}}$.
\item Scattering:
\[
\langle f|iT|i\rangle=i(2\pi)^4\delta^{(4)}(p_f-p_i)\mathcal{M}.
\]
\item Decays and cross sections computed from $|\mathcal{M}|^2$ and phase space.
\item Vacuum bubbles cancel in normalized correlators.
\item LSZ relates $S$-matrix elements to amputated time-ordered Green functions.
\end{itemize}
